\section{The Anchor Point Method (the centerpiece)}
\label{sec:anchor}

This is the \textbf{centerpiece} of the project --- the clever idea that makes
the whole thing fast \emph{and} trustworthy. It lives in
\file{lunar/equilibrium.py}. Read that file's docstring; it's excellent. This
section builds on it.

\subsection{The problem: the deep Moon settles agonizingly slowly}

We want the \textbf{periodic steady state}: the repeating day-after-day
temperature pattern the soil settles into after a long time, when every lunar day
looks like the last. To get there the obvious approach is ``spin-up'' --- just
press play and simulate many lunar days until things stop changing.

The catch: the \textbf{deep} column relaxes on the \emph{diffusion timescale of
the whole 5-metre column}, which is roughly \textbf{1000 lunar days}. Simulating
1000 lunations per run, for the ${\sim}300$ runs the retrieval needs, is
hopeless.

Worse, if you stop early (say after 30 lunations), the deep temperature still
\textbf{secretly remembers your starting guess}. And the deep temperature is
exactly what the Apollo sensors measure. So your ``measurement'' of \Kd{} would
be contaminated by an arbitrary guess you typed in. This was a real, critical bug
caught in the project's own audit --- \textbf{flag F1} in
\file{docs/FLAG\_REPORT.md}. The original code used hard-coded starting
temperatures (\SI{250}{\kelvin} for A15, \SI{255}{\kelvin} for A17), and a
$\pm\SI{5}{\kelvin}$ change in that guess swung the retrieved \Kdstar{} by a
factor of ${\sim}6$ --- \emph{larger than the very inter-site contrast the paper
reports!} That had to be fixed before any result could be believed.

\subsection{The key insight: two clocks}
\label{sec:twoclocks}

The fix comes from realizing the soil has \textbf{two very different ``clocks'':}
\begin{itemize}[leftmargin=1.4em,itemsep=3pt]
  \item \textbf{The shallow skin (top ${\sim}$half metre) is FAST.} It settles
    into its repeating day/night rhythm within a handful of lunar days, because
    it is thin and driven hard by the surface.
  \item \textbf{The deep column is SLOW} to settle by brute force --- \emph{but}
    in steady state it must obey one dead-simple rule: the same steady interior
    heat \Qb{} flows through every layer. That single rule,
    $\mathrm{d}\langle T\rangle/\mathrm{d}z = Q_b/K$, \textbf{fixes the entire
    deep temperature shape instantly} --- no waiting required. ($\langle T\rangle$
    means the cycle-averaged temperature; the angle brackets denote ``average
    over one day.'')
\end{itemize}
So instead of waiting 1000 lunations for the deep column, we \textbf{calculate}
it from the steady-flux rule, and only ever simulate the fast skin.

\subsubsection{Where $\mathrm{d}\langle T\rangle/\mathrm{d}z = Q_b/K$ comes from}

Same physics as \autoref{sec:boundaries} (geothermal gradient), now as the
anchor shortcut:
\begin{enumerate}[leftmargin=1.5em,itemsep=2pt]
  \item Time-average the heat equation over one lunation.
  \item In periodic steady state: $\partial\langle T\rangle/\partial t=0$.
  \item Equation collapses to $\partial[K\,\partial\langle T\rangle/\partial z]/\partial z=0$.
  \item So $K\,\mathrm{d}\langle T\rangle/\mathrm{d}z=Q_b$ everywhere below the skin.
  \item Integrate downward from the anchor: \code{\_reconstruct\_subskin} in
    \file{lunar/equilibrium.py}.
\end{enumerate}

\subsubsection{The rectified flux $u_{\mathrm{rect}}$}

Inside the skin, $\langle K\,\partial T/\partial z\rangle \neq K(\langle T\rangle)\,
\mathrm{d}\langle T\rangle/\mathrm{d}z$. The extra \textbf{rectification} term
$u_{\mathrm{rect}}$ is measured from the resolved cycle (\code{\_rectified\_flux})
and subtracted:
$\mathrm{d}\langle T\rangle/\mathrm{d}z = (Q_b - u_{\mathrm{rect}})/K$.
Negligible below \SI{1}{m}; a few percent near \SI{0.3}{m}--\SI{0.5}{m} --- why
the anchor sits at \SI{0.55}{m}.

\begin{figure}[h]
\centering
\begin{tikzpicture}[x=1cm,y=1cm,font=\small]
  % skin box
  \fill[coral!12] (0,0) rectangle (5,-1.6);
  \draw[coral,line width=0.9pt] (0,0) rectangle (5,-1.6);
  \node[coral,font=\bfseries] at (2.5,-0.45) {FAST skin (top ${\sim}0.5$ m)};
  \node[char,font=\scriptsize,align=center] at (2.5,-1.15)
    {settles in a handful of lunations\\ {\itshape step 1: simulate it fully}};
  % deep box
  \fill[teal!12] (0,-1.8) rectangle (5,-6.2);
  \draw[teal,line width=0.9pt] (0,-1.8) rectangle (5,-6.2);
  \node[teal,font=\bfseries] at (2.5,-2.25) {SLOW deep column};
  \node[char,font=\scriptsize,align=center] at (2.5,-3.1)
    {brute force ${\approx}1000$ lunations\\ {\itshape step 2: don't wait --- compute it}};
  % flux rule
  \node[teal,font=\small] at (2.5,-4.4) {$\dfrac{\mathrm{d}\langle T\rangle}{\mathrm{d}z}=\dfrac{Q_b-u_{\mathrm{rect}}}{K}$};
  \draw[-{Latex[length=2.5mm]},teal,line width=1pt] (2.5,-4.95) -- (2.5,-6.0)
     node[midway,right,font=\scriptsize] {integrate downward};
  % anchor
  \draw[dashed,plum,line width=1pt] (-0.4,-1.8) -- (5.4,-1.8);
  \node[plum,font=\bfseries\scriptsize,anchor=west] at (5.5,-1.8) {anchor depth (0.55 m)};
  % basal flux arrow
  \draw[-{Latex[length=3mm]},forest,line width=1.2pt] (2.5,-6.9) -- (2.5,-6.25)
     node[below,font=\scriptsize,yshift=-8pt,forest] {basal heat flux $Q_b$ (from the deep interior)};
  % outer loop arrows on the side
  \draw[-{Latex[length=2mm]},dim,line width=0.8pt]
     (5.9,-3.5) to[out=60,in=-60] (5.9,-1.0);
  \node[dim,font=\scriptsize,rotate=90,anchor=south] at (6.25,-2.3) {repeat 3--5$\times$};
\end{tikzpicture}
\caption[The two-clock Anchor Point Method]{The Anchor Point Method as ``two
clocks.'' \textbf{What to notice:} the fast skin is simulated in full (step~1);
the slow deep column is \emph{reconstructed} from the steady-flux rule
(step~2) rather than waited out, anchored to the cycle-mean temperature read at
$0.55$~m. The two steps are repeated a few times until they agree.
\textbf{Why it matters:} this turns an impossible ${\sim}1000$-lunation wait into
a handful of cheap iterations, while removing the starting-guess bias.}
\label{fig:twoclocks}
\end{figure}

\subsection{The method: alternate ``settle the skin'' and ``rebuild the deep''}

The algorithm (the ``outer iteration'') alternates two cheap steps until they
agree:
\begin{enumerate}[leftmargin=1.5em,itemsep=4pt]
  \item \textbf{Settle the skin.} Run the full solver for just a few lunar days
    (\code{n\_inner = 12}) starting from the current profile. This locks the fast
    skin into rhythm against the current deep column.
  \item \textbf{Rebuild the deep.} Read the cycle-mean temperature at one trusted
    \textbf{anchor depth} just below the skin (\code{z\_anchor = 0.55} m), then
    integrate the steady-flux rule
    $\mathrm{d}\langle T\rangle/\mathrm{d}z = (Q_b - u_{\mathrm{rect}})/K$
    \emph{downward} from there to reconstruct the entire deep profile exactly.
\end{enumerate}
Repeat steps 1--2 a few times (typically 3--5). The result ``locks on'' to the
true periodic steady state. The actual loop is \code{solve\_periodic\_equilibrium}
in \file{lunar/equilibrium.py}; the deep rebuild is \code{\_reconstruct\_subskin}.

Two refinements worth knowing:
\begin{itemize}[leftmargin=1.4em,itemsep=3pt]
  \item \textbf{$u_{\mathrm{rect}}$ --- the rectified flux.} Inside the skin, the
    daily wobble plus the fact that conductivity depends on temperature creates a
    small extra \emph{day--night-averaged} heat term (an ``eddy'' or
    ``rectification'' term). The code measures it exactly from the resolved cycle
    (\code{\_rectified\_flux}) and subtracts it, so the deep rebuild stays
    accurate. It is negligible below ${\sim}1$~m but matters a few percent right
    at the anchor --- which is precisely why the anchor is placed at
    $0.55$~m, \emph{below} the worst of it.
  \item \textbf{Two-stage anchor schedule.} The code first anchors high
    ($0.25$~m, relaxes fast) to cheaply kill most of the starting-guess error,
    then drops to $0.55$~m for the accurate finish. See the \code{stages} tuple
    in \code{solve\_periodic\_equilibrium}.
\end{itemize}

\subsection{Why it's both fast \emph{and} unbiased}

\begin{itemize}[leftmargin=1.4em,itemsep=3pt]
  \item \textbf{Fast:} it only ever simulates the quick skin (a few lunations per
    outer step, ${\sim}36$--$60$ lunations total) instead of ${\sim}1000$.
    Comparable cost to the old broken spin-up, but correct.
  \item \textbf{Unbiased (this is the crucial part):} the final answer \textbf{no
    longer depends on the starting guess}. The skin is forced periodic by step~1;
    the deep is forced to satisfy steady flux by step~2. The only fixed point of
    that two-step map is the \emph{true} equilibrium. The starting guess just
    seeds the first iterate and gets erased.
\end{itemize}

This is \textbf{certified, not assumed.} The regression test
\code{tests/test\_equilibrium.py::test\_guess\_independence} literally runs the
solver from two different starting guesses (\SI{242}{\kelvin} and
\SI{258}{\kelvin}) and asserts the converged profiles agree:

\begin{lstlisting}[caption={The guess-independence regression test.}]
def test_guess_independence():
    """Converged profiles from 242 K and 258 K guesses must agree."""
    grid, t, insol, k_func, cp_func = _setup()
    profiles = []
    for guess in (242.0, 258.0):
        eq = solve_periodic_equilibrium(
            grid=grid, t=t, insolation=insol, albedo=0.131,
            emissivity=0.95, Q_b=0.021, K_func=k_func, cp_func=cp_func,
            T_guess=guess)
        assert eq.converged
        profiles.append(eq.T_mean)
    assert np.max(np.abs(profiles[0] - profiles[1])) < 0.1
\end{lstlisting}

In the full production setting the guess-dependence is $\le \SI{0.023}{\kelvin}$
(vs the old \SI{1.6}{}--\SI{4.5}{\kelvin}). That is what ``F1 RESOLVED'' at the
top of \file{docs/FLAG\_REPORT.md} means, and it's why the retrieved \Kdstar{}
values can be trusted.

\subsection{``Flux closure'' --- how the code proves it worked}

How do you know the deep profile really satisfies steady heat flow? You check
whether the cycle-mean conductive flux equals \Qb{} at every depth. The maximum
relative deviation is the \textbf{flux closure} number (\code{flux\_closure} in
\code{EquilibriumResult}). The retrieval warns if closure is worse than $3\%$. It
is deliberately measured \textbf{below the diurnal skin} (where the rule actually
applies and where every Apollo sensor sits), not at the noisy surface. The test
\code{test\_flux\_closure} checks this. The
\code{make\_equilibrium\_certification.py} figure script visualizes the whole
certification.
